\documentclass{article}
\usepackage[utf8]{inputenc}
\usepackage{amsmath}
\usepackage{amsfonts}
\usepackage{amssymb}
\usepackage{geometry}
\geometry{a4paper, margin=1in}

\title{MAC0344 Arquitetura de Computadores \\ Lista de Exercícios No. 3}
\author{Leonardo Heidi Almeida Murakami \\ NUSP: 11260186}
\date{\today}

\begin{document}
\maketitle

\section*{Exercício 1}
\subsection*{a}
Qualquer apartamento no centro tem aluguel mais baixo do que alguns apartamentos em Pinheiros
\begin{equation}
    \forall x((\text{Apto}(x)\wedge\text{Em}(x,\text{centro}))) \rightarrow \exists y ((\text{Apto}(y)\wedge\text{Em}(y, \text{pinheiros}))) \rightarrow (\text{Aluguel}(x) < \text{Aluguel}(y))
\end{equation}
Não é uma boa tradução, uma tradução correta seria
\begin{equation}
    \forall x((\text{Apto}(x)\wedge\text{Em}(x,\text{centro}))) \rightarrow \exists y((\text{Apto}(y) \wedge \text{Em}(y, \text{pinheiros})) \wedge \text{Aluguel}(x) < \text{Aluguel}(y))
\end{equation}


\subsection*{b}
"Há exatamente um apartamento em Pinheiros com aluguel abaixo de R\$1000."
\begin{equation}
    \exists x(\text{Apto}(x) \wedge \text{Em}(x, \text{pinheiros}) \wedge
    \forall y((\text{Apto}(y) \wedge \text{Em}(y, \text{pinheiros}) \wedge (\text{Aluguel}(y) < 1000)) \rightarrow (y = x)))
\end{equation}

\section{Exercicio 2}
\begin{itemize}
    \item \textbf{esp} -- Constante, naipe de espadas.
    \item \textbf{m1} -- Constante, uma mão particular.
    \item $n(x, y)$ -- O naipe da carta $x$ é $y$.
    \item $i(x, y)$ -- A carta $x$ está na mão $y$.
    \item $f(x)$ -- A mão $x$ é um flush.
\end{itemize}

\begin{enumerate}
    \item[(a)] \textbf{Traduza para a lógica de primeira ordem:}
    \begin{enumerate}
        \item[(i)] Se todas as cartas de uma mão são do mesmo naipe, a mão é um flush.
        \item[(ii)] Toda carta de $m_1$ é de espadas.
    \end{enumerate}
    \item[(b)] \textbf{Use resolução para provar que $m_1$ é um flush a partir de (i) e (ii).}
\end{enumerate}

\subsection*{a}
\begin{itemize}
    \item Se todas as cartas de uma mão são do mesmo naipe, a mão é um flush.
    \item $\forall x \left( \left( \exists y \; \forall z \; (i(z, x) \rightarrow n(z, y)) \right) \rightarrow f(x) \right)$
\end{itemize}
\end{document}